% XeLaTeX twice. This file is an editable, watermark-free regression specimen.
\documentclass[11pt,compress]{beamer}
\usepackage{ruc-research}
\title[无水印模板测试]{中文科研组会模板}
\subtitle{Chinese text, equations, tables and figures}
\author[模板测试]{模板功能测试}
\institute[科研组会]{可复用的 RUC Beamer 工程}
\date{2026年11月15日}
\begin{document}
\begin{frame}[plain,noframenumbering]\titlepage\end{frame}
\section{正文}
\begin{frame}{公式与中英文混排}
中文、English、$\alpha$、$\beta$、$R^2$ 与百分号 5\%。
\begin{equation}\label{eq:example}
 R_{i,t}=\alpha_i+\beta_i R_{m,t}+\varepsilon_{i,t}.
\end{equation}
式\eqref{eq:example}中的 $\beta_i$ 表示市场暴露。
\source{排版测试示例；不是论文估计结果。}
\end{frame}
\begin{frame}{三线表、编号与引用}
\begin{table}\centering
\caption{示例数据}\label{tab:test}
\begin{tabular}{lrr}\toprule
项目 & 样本 A & 样本 B\\\midrule
演示值 & 1.23 & 4.56\\
样本数 & 100 & 120\\\bottomrule
\end{tabular}
\end{table}
表\ref{tab:test}仅测试格式。
\end{frame}
\begin{frame}{图形、编号与可编辑对象}
\begin{figure}\centering
\begin{tikzpicture}[>=Stealth]
\node[draw=tsinghua,minimum width=2.6cm,minimum height=1cm] (a) {研究问题};
\node[draw=tsinghua,right=1.2cm of a,minimum width=2.6cm,minimum height=1cm] (b) {可核查证据};
\draw[->,thick,tsinghua] (a)--(b);
\end{tikzpicture}
\caption{流程示意图}\label{fig:test}
\end{figure}
图\ref{fig:test}为 TikZ 矢量图。背景没有校徽。
\end{frame}
\appendix
\begin{frame}{附录独立计数}
正文总页数为 3；本页应显示 A1/1，封面不计数。
\end{frame}
\end{document}
